\section{Introduction}\label{sec:intro}

Start with an $h\times m$ grid and black out some of its points.
The game master chooses a rectangle (tilted ones are allowed) and tells
you which of its corners are white. Can you recover all four corners?
Kagey~\cite{kageyG4G14,kagey001} asks how many points can be blacked out
while keeping the answer unique for every rectangle.
Such a set is a \emph{valid blackout} (Definition~\ref{def:valid}).
One round is shown in Figure~\ref{fig:game}.

\begin{figure}[t]
\centering
\resizebox{\textwidth}{!}{\begin{tikzpicture}[x=0.9cm,y=0.9cm,baseline=(current bounding box.north)]
  \begin{scope}[xshift=0.000cm]
    \draw[black,line width=0.36pt,dash pattern=on 1.35pt off 1.35pt] (1,-1) -- (4,-1) (1,-2) -- (4,-2) (1,-3) -- (4,-3) (1,-1) -- (1,-3) (2,-1) -- (2,-3) (3,-1) -- (3,-3) (4,-1) -- (4,-3);
    \filldraw[fill=white,draw=black,line width=0.54pt] (1,-1) circle (0.3);
    \filldraw[fill=white,draw=black,line width=0.54pt] (2,-1) circle (0.3);
    \filldraw[fill=white,draw=black,line width=0.54pt] (3,-1) circle (0.3);
    \filldraw[fill=white,draw=black,line width=0.54pt] (4,-1) circle (0.3);
    \filldraw[fill=white,draw=black,line width=0.54pt] (1,-2) circle (0.3);
    \filldraw[fill=white,draw=black,line width=0.54pt] (2,-2) circle (0.3);
    \filldraw[fill=kyellow,draw=black,line width=0.54pt] (3,-2) circle (0.3);
    \filldraw[fill=white,draw=black,line width=0.54pt] (4,-2) circle (0.3);
    \filldraw[fill=white,draw=black,line width=0.54pt] (1,-3) circle (0.3);
    \filldraw[fill=white,draw=black,line width=0.54pt] (2,-3) circle (0.3);
    \filldraw[fill=white,draw=black,line width=0.54pt] (3,-3) circle (0.3);
    \filldraw[fill=white,draw=black,line width=0.54pt] (4,-3) circle (0.3);
    \node[lab] at (1,-0.45) {1};
    \node[lab] at (2,-0.45) {2};
    \node[lab] at (3,-0.45) {3};
    \node[lab] at (4,-0.45) {4};
    \node[lab] at (0.45,-1) {1};
    \node[lab] at (0.45,-2) {2};
    \node[lab] at (0.45,-3) {3};
    \node[capt,anchor=north] at (2.50,-3.45) {(a) the point $(3,2)$};
  \end{scope}
  \begin{scope}[xshift=4.680cm]
    \draw[black,line width=0.36pt,dash pattern=on 1.35pt off 1.35pt] (1,-1) -- (4,-1) (1,-2) -- (4,-2) (1,-3) -- (4,-3) (1,-1) -- (1,-3) (2,-1) -- (2,-3) (3,-1) -- (3,-3) (4,-1) -- (4,-3);
    \filldraw[fill=kgray,draw=black,line width=0.54pt] (1,-1) circle (0.3);
    \draw[black,line width=0.54pt] (0.788,-0.788) -- (1.212,-1.212) (0.788,-1.212) -- (1.212,-0.788);
    \filldraw[fill=kgray,draw=black,line width=0.54pt] (2,-1) circle (0.3);
    \draw[black,line width=0.54pt] (1.788,-0.788) -- (2.212,-1.212) (1.788,-1.212) -- (2.212,-0.788);
    \filldraw[fill=kgray,draw=black,line width=0.54pt] (3,-1) circle (0.3);
    \draw[black,line width=0.54pt] (2.788,-0.788) -- (3.212,-1.212) (2.788,-1.212) -- (3.212,-0.788);
    \filldraw[fill=kgray,draw=black,line width=0.54pt] (4,-1) circle (0.3);
    \draw[black,line width=0.54pt] (3.788,-0.788) -- (4.212,-1.212) (3.788,-1.212) -- (4.212,-0.788);
    \filldraw[fill=kgray,draw=black,line width=0.54pt] (1,-2) circle (0.3);
    \draw[black,line width=0.54pt] (0.788,-1.788) -- (1.212,-2.212) (0.788,-2.212) -- (1.212,-1.788);
    \filldraw[fill=white,draw=black,line width=0.54pt] (2,-2) circle (0.3);
    \filldraw[fill=white,draw=black,line width=0.54pt] (3,-2) circle (0.3);
    \filldraw[fill=white,draw=black,line width=0.54pt] (4,-2) circle (0.3);
    \filldraw[fill=kgray,draw=black,line width=0.54pt] (1,-3) circle (0.3);
    \draw[black,line width=0.54pt] (0.788,-2.788) -- (1.212,-3.212) (0.788,-3.212) -- (1.212,-2.788);
    \filldraw[fill=white,draw=black,line width=0.54pt] (2,-3) circle (0.3);
    \filldraw[fill=white,draw=black,line width=0.54pt] (3,-3) circle (0.3);
    \filldraw[fill=white,draw=black,line width=0.54pt] (4,-3) circle (0.3);
    \node[lab] at (1,-0.45) {1};
    \node[lab] at (2,-0.45) {2};
    \node[lab] at (3,-0.45) {3};
    \node[lab] at (4,-0.45) {4};
    \node[lab] at (0.45,-1) {1};
    \node[lab] at (0.45,-2) {2};
    \node[lab] at (0.45,-3) {3};
    \node[capt,anchor=north] at (2.50,-3.45) {(b) our blackout};
  \end{scope}
  \begin{scope}[xshift=9.360cm]
    \draw[black,line width=0.36pt,dash pattern=on 1.35pt off 1.35pt] (1,-1) -- (4,-1) (1,-2) -- (4,-2) (1,-3) -- (4,-3) (1,-1) -- (1,-3) (2,-1) -- (2,-3) (3,-1) -- (3,-3) (4,-1) -- (4,-3);
    \draw[kred,line width=2.55pt,line join=round] (2,-2) -- (3,-1) -- (4,-2) -- (3,-3) -- cycle;
    \filldraw[fill=kgray,draw=black,line width=0.54pt] (1,-1) circle (0.3);
    \draw[black,line width=0.54pt] (0.788,-0.788) -- (1.212,-1.212) (0.788,-1.212) -- (1.212,-0.788);
    \filldraw[fill=kgray,draw=black,line width=0.54pt] (2,-1) circle (0.3);
    \draw[black,line width=0.54pt] (1.788,-0.788) -- (2.212,-1.212) (1.788,-1.212) -- (2.212,-0.788);
    \filldraw[fill=kgray,draw=black,line width=0.54pt] (3,-1) circle (0.3);
    \draw[black,line width=0.54pt] (2.788,-0.788) -- (3.212,-1.212) (2.788,-1.212) -- (3.212,-0.788);
    \filldraw[fill=kgray,draw=black,line width=0.54pt] (4,-1) circle (0.3);
    \draw[black,line width=0.54pt] (3.788,-0.788) -- (4.212,-1.212) (3.788,-1.212) -- (4.212,-0.788);
    \filldraw[fill=kgray,draw=black,line width=0.54pt] (1,-2) circle (0.3);
    \draw[black,line width=0.54pt] (0.788,-1.788) -- (1.212,-2.212) (0.788,-2.212) -- (1.212,-1.788);
    \filldraw[fill=white,draw=black,line width=0.54pt] (2,-2) circle (0.3);
    \filldraw[fill=white,draw=black,line width=0.54pt] (3,-2) circle (0.3);
    \filldraw[fill=white,draw=black,line width=0.54pt] (4,-2) circle (0.3);
    \filldraw[fill=kgray,draw=black,line width=0.54pt] (1,-3) circle (0.3);
    \draw[black,line width=0.54pt] (0.788,-2.788) -- (1.212,-3.212) (0.788,-3.212) -- (1.212,-2.788);
    \filldraw[fill=white,draw=black,line width=0.54pt] (2,-3) circle (0.3);
    \filldraw[fill=white,draw=black,line width=0.54pt] (3,-3) circle (0.3);
    \filldraw[fill=white,draw=black,line width=0.54pt] (4,-3) circle (0.3);
    \node[lab] at (1,-0.45) {1};
    \node[lab] at (2,-0.45) {2};
    \node[lab] at (3,-0.45) {3};
    \node[lab] at (4,-0.45) {4};
    \node[lab] at (0.45,-1) {1};
    \node[lab] at (0.45,-2) {2};
    \node[lab] at (0.45,-3) {3};
    \node[capt,anchor=north] at (2.50,-3.45) {(c) the rectangle};
  \end{scope}
  \begin{scope}[xshift=14.040cm]
    \draw[black,line width=0.36pt,dash pattern=on 1.35pt off 1.35pt] (1,-1) -- (4,-1) (1,-2) -- (4,-2) (1,-3) -- (4,-3) (1,-1) -- (1,-3) (2,-1) -- (2,-3) (3,-1) -- (3,-3) (4,-1) -- (4,-3);
    \filldraw[fill=kyellow,draw=black,line width=0.54pt] (2,-2) circle (0.3);
    \filldraw[fill=kyellow,draw=black,line width=0.54pt] (4,-2) circle (0.3);
    \filldraw[fill=kyellow,draw=black,line width=0.54pt] (3,-3) circle (0.3);
    \node[lab] at (1,-0.45) {1};
    \node[lab] at (2,-0.45) {2};
    \node[lab] at (3,-0.45) {3};
    \node[lab] at (4,-0.45) {4};
    \node[lab] at (0.45,-1) {1};
    \node[lab] at (0.45,-2) {2};
    \node[lab] at (0.45,-3) {3};
    \node[capt,anchor=north] at (2.50,-3.45) {(d) what we are shown};
  \end{scope}
\end{tikzpicture}
}
\caption{One round on a largest valid blackout of the $3\times4$ grid
(Section~\ref{sec:three}). Points are (column, row). Gray crossed circles are black. Panel (c) shows the red
rectangle, and (d) its yellow corners. The three yellow corners fix the fourth by
Lemma~\ref{lem:threecorners}.}
\label{fig:game}
\end{figure}

Write $f(h,m)$ for the most points we can black out, and $a_h(m)$
for the number of ways to attain it. The conjectures recorded in our
OEIS entries \seqnum{A397315} and \seqnum{A399580}~\cite{oeis} were
$f(3,m)=m+2$ and $f(h,m)=h+m$ for $\min(h,m)\ge4$, with exceptions
$f(4,4)=7$ and $f(4,5)=8$. The first was proved in Lean by the Omega
Institute, whose pull request was merged on 15 September
2026~\cite{trureturing2026}. Our independent proof combines the row
bound (Theorem~\ref{thm:rowbound}) and the boundary cross
(Lemma~\ref{lem:cross}). The second guess breaks at $5\times8$
(Theorem~\ref{thm:fiverow}). Throughout, put $q_h=\lfloor h^2/4\rfloor$.
Computer-assisted claims are marked, with evidence in
Table~\ref{tab:evidence}.
% TODO(author): add the date of A397315 after its A-number (item S6).
% TODO(author): A-numbers for rows 3, 4 of A399580.

The main results follow, with proof locations in parentheses.

\begin{introthm}[narrow grids]\label{thm:A}
\leavevmode
\begin{enumerate}[(a)]
\item We have $f(h,m)\ge h+m-1$, and $f(h,m)\le m+q_h$ if $m\ge3$. \textup{(Lemma~\ref{lem:cross}
and Theorem~\ref{thm:rowbound}, Section~\ref{sec:graph}.)}
\item We have $f(2,m)=m+1$ and $f(3,m)=m+2$ for $m\ge3$, $f(4,m)=m+4$ for $m\ge6$,
and $f(5,5)=10$. \textup{(Theorem~\ref{thm:strip}, Section~\ref{sec:three},
Theorem~\ref{thm:fourvalue}, and Theorem~\ref{thm:fivefive}(a).)}
\item (Computer-assisted, with checked certificates for $f(5,6)$ and $f(5,7)$, and the last
claim checked in Lean.) We have $f(4,4)=7$, $f(4,5)=8$,
$f(5,6)=11$, $f(5,7)=12$, and $f(5,m)=m+6$ for $m\ge8$. \textup{(Corollary~\ref{cor:fourboundary}
and Theorem~\ref{thm:fiverow}(a), (b).)}
\item (Computer-assisted for $m\ge9$, with checked certificates for $m\le14$, and checked in
Lean for $m\ge15$.) We have $f(6,m)\le m+8$ for $8\le m\le14$, and
$f(6,m)=m+9$ for $m\ge15$. \textup{(Theorem~\ref{thm:sixrow}.)}
\item (Computer-assisted, with checked certificates.) We have $f(6,6)=12$,
$f(6,7)=13$, $f(7,7)=14$, and $f(6,8)=15$. \textup{(Theorem~\ref{thm:refute}.)}
\end{enumerate}
\end{introthm}

For (c) and (d), the collar theorem reduces each infinite family to
one validity check (Theorem~\ref{thm:collar}). Lemma~\ref{lem:enc}
justifies the SAT bounds.

\begin{introthm}[three-row optima]\label{thm:B}
\leavevmode
\begin{enumerate}[(a)]
\item For $m\ge4$ we have $a_3(m)=4(m-1)(3m+2)\,3^{m-5}+E_m$, where $E_m>0$ counts the
optima with a column black in rows $1$ and $3$ (\seqnum{A398448}), and
$E_m=O(m^2(1+\sqrt3)^m)$. So $a_3(m)\sim\frac4{81}m^23^m$, and a fraction tending to $4/27$
of the $m^23^{m-1}$ axis-valid candidates is valid. The series $\sum_{m\ge3}a_3(m)z^m$ is
rational, and its denominator divides $(1-3z)^3(1-2z-2z^2)^3(1+2z^2)^3$.
\textup{(Theorem~\ref{thm:threecount}, Proposition~\ref{prop:threeclass}(c), the last claims of
Theorem~\ref{thm:threeclosed}, and Theorem~\ref{thm:threegf}(a), Section~\ref{sec:three}.)}
\item (Computer-assisted, exact computer algebra.) This is the reduced denominator,
and $E_m\sim\frac{19+11\sqrt3}{432}m^2(1+\sqrt3)^m$. \textup{(Theorem~\ref{thm:threegf}(b),
proved in Section~\ref{app:threegf}, and Theorem~\ref{thm:threeclosed}.)}
\end{enumerate}
\end{introthm}

\begin{introthm}[four rows]\label{thm:C}
If $m\ge4$, no valid blackout of the $4\times m$ grid hides all corners of a
rectangle \textup{(Theorem~\ref{thm:hidden})}. (Computer-assisted.) The series
$\sum_{m\ge6}a_4(m)z^m$ is rational with reduced denominator of degree $420$, and
$a_4(m)\sim c\,m^4\rho_4^m$ with $c\approx8.1611\times10^{-5}$ and $\rho_4=3.35199\ldots$, an
algebraic number of degree $20$ \textup{(Proposition~\ref{prop:fourautomaton},
Theorem~\ref{thm:fourgfexact}, and Corollary~\ref{cor:fourgrowth}, Section~\ref{sec:four})}.
\end{introthm}

\begin{introthm}[every fixed height]\label{thm:E}
Fix $h\ge1$ and write a blackout as the word $B_1\cdots B_m$ of the black rows of its
columns. The valid words form a regular language. If $h\ge2$, then $f(h,m)-m$ is
eventually periodic in $m$, and $\sum_{m\ge0}a_h(m)z^m$ is rational.
\textup{(Theorems~\ref{thm:regular} and~\ref{thm:periodic}, Section~\ref{sec:height}.)}
\end{introthm}

Corollary~\ref{cor:oeisconj} proves the conjecture in \seqnum{A338885} that its row $n$ has
largest entry \seqnum{A033638}$(n)$.

\paragraph{Prior work.} Choosing white points to distinguish rectangles is
a test-cover problem~\cite{bellitto2018,foucaud2015,katona1966,moretshapiro1985}
(Corollary~\ref{cor:hitting}). If we only forbid fully black axis-aligned
rectangles, we get the Zarankiewicz problem~\cite{reiman1958}.
To distinguish all axis-aligned rectangles, the row-column graph must
avoid both four- and six-cycles (Lemma~\ref{lem:girth}). Such graphs,
including unbalanced parts, are studied by de Caen and
Sz\'ekely~\cite{decaenszekely1992,decaenszekely1997}.
Neuwirth's edge bounds~\cite[Theorem~1(ii) and Proposition~17]{neuwirth2001}
already give our row bound in graph form. Our Mantel proof~\cite{mantel1907} also classifies and counts
the equality cases.
The tilted rectangles require additional arguments.
For fixed height, our column method uses transfer
matrices~\cite{alanko2011,euler2005,mertens2024},
\cite[Proposition~V.6]{flajolet2009}.
Theorem~1(3) of Guilhot, Little, and Parkinson~\cite{glp2025} also
yields Theorem~\ref{thm:E} once regularity is proved.

\paragraph{Outline.} Sections~\ref{sec:conv} to~\ref{sec:graph} set up
the problem and its graph bounds. Sections~\ref{sec:three}
to~\ref{sec:height} handle narrow grids and fixed height.
Section~\ref{sec:verify} gives computational evidence and encoding
proofs, Section~\ref{sec:open} the open questions, and
Appendix~\ref{app:square} the square-grid bound.
